\chapter{Building a Real Screen}
\label{ch:real-screen}
\chapterquote{The value of a framework is not visible in the first field; it is visible when the tenth concern arrives and the screen remains legible.}{The screen-building rule}

\begin{chaptermap}
Core question & How do the book's pieces combine into one maintainable Swing workspace that a developer can inspect, test, extend, and screenshot?\\
Ingredients & Records, generated forms, generated table descriptors, commands, MigLayout, FlatLaf, validation, table state, background loading, busy UI, lifecycle scopes, resources, help, and tests.\\
Plain Swing baseline & A handwritten screen can work, but row data, selection, validation, command state, table state, task state, and cleanup need explicit owners.\\
Corusco connection & Corusco supplies stable facts, observable state, generated companions, binding scopes, command descriptors, table descriptors, problem models, and task boundaries while layout and workflow remain readable application code.\\
Companion example & \bookAppExample.\\
\end{chaptermap}

\section{The target screen}

The running screen in this chapter is a customer workspace. It has a searchable
customer table, a detail form, commands, status text, validation feedback,
saved table layout, background refresh, busy indication, and tests. This is
not a large application, but it is large enough to expose the important
interactions. A user will search, sort, select, edit, save, cancel, refresh,
and wait. A developer will add a column, rename a label, change a validation
rule, persist table state, route F1 help, and write a screenshot test. The
screen is deliberately ordinary because ordinary business screens are where
most Swing maintainability is won or lost.

The companion \bookAppExample{} is a compact executable shell for screenshots:
FlatLaf-oriented construction, MigLayout geometry, a customer table, a detail
panel, deterministic sample data, and a busy-detail variant. The full chapter
describes the production shape that this shell is meant to grow into. That
distinction matters. A book example should be small enough to read. The
architecture should be strong enough for the tenth concern.

\begin{figure}[htbp]
\centering
\begin{minipage}[t]{0.48\linewidth}
\includegraphics[width=\linewidth]{\bookAppFigure}
\end{minipage}\hfill
\begin{minipage}[t]{0.48\linewidth}
\includegraphics[width=\linewidth]{\bookAppBusyFigure}
\end{minipage}
\caption{The real-screen companion appears in normal and busy-detail states. The screenshots keep the chapter grounded in an actual FlatLaf and MigLayout Swing surface.}
\label{fig:bookapp-workspace-screens}
\end{figure}

Figure~\ref{fig:bookapp-workspace-screens} is a visual contract. The table is
not a placeholder for a future web grid. The detail form is not a mockup. The
screen is a real Swing surface whose layout, components, and state can be
constructed on the EDT, tested, and reviewed as a deterministic screenshot. The chapter will
speak about generated form models, command descriptors, task services, and
table-state controllers, but the reader should keep the visible workspace in
mind. The abstractions are justified only if this kind of screen remains
clearer after they are added.

The screen is also intentionally modest in style. It is not a marketing page
and not a visual-effects demo. A user must scan rows, compare fields, type
edits, and read validation. MigLayout gives the main regions their proportions.
FlatLaf gives modern widgets and density. Corusco gives stable identities and
behavioral contracts. The application still decides what the user is trying to
do and how the workflow should feel.

\termdef{Workspace}{in this chapter} means a Swing view that contains several
cooperating areas rather than a single modal form. A workspace often has a
table or tree on one side, detail editors on another side, commands around the
edges, status text, background activity, and local navigation. It may live in a
frame, internal frame, tab, or larger shell. The exact container is less
important than the fact that several concerns meet on one screen.

\termdef{Concern ownership}{in this chapter} means that each durable fact has
a named owner. Rows belong to a row source. Selection belongs to a selected-row
value or selection model. Editable detail belongs to a form model. Commands
belong to command objects. Long-running work belongs to task handles and
lifecycles. Table layout belongs to table state. Visible text belongs to
resources. A listener may translate an event, but it should rarely be the only
place where a durable fact exists.

The chapter's main question is not "which helper class should I import?" It is
"where does this fact live after the screen becomes real?" If the answer is
"in the current text of a button", "in column index two", "in the last
listener that ran", or "in a background callback that may arrive later", the
screen will become difficult to maintain. If the answer is an observable
value, generated key, descriptor, command, scope, table-state store, resource
key, or task handle, the screen has a structure that can be reviewed.

This question should be asked while the screen is still small. A three-field
form can survive a listener per field. A five-row table can survive an array
model. A single Refresh button can survive a direct service call. The danger
is that working code becomes trusted structure. The next feature then attaches
itself to the nearest listener because that seems cheaper than naming the
missing owner. The customer workspace is built as if it will grow, because
most successful desktop screens do grow.

There is also a reading contract. A future maintainer should be able to open
the workspace and find the main facts quickly. The row type should be visible.
The generated table descriptor should be visible. The form model should be
visible. The command set should be visible. The task boundary should be
visible. The view layout should be visible. The visual test state should be
visible. If understanding the screen requires following anonymous listeners
through component construction order, the code may work but the design
standard has not been met.

\section{Stable facts before components}

Begin with the facts that should survive layout and localization: customer row
identity, table id, column ids, form field ids, command ids, resource ids, help
topics, validation problem codes, and component keys for tests. These facts
should not be discovered from button text, table headers, or current component
positions. They belong in records, annotations, generated companions, or
handwritten descriptors. The component tree comes later.

Stable facts are cheap to define early and expensive to recover late. Once a
table is full of localized headers and integer indexes, it is hard to recover
a durable column identity. Once tests locate a Save button by text, changing
the resource bundle becomes risky. Once a validation message is the only thing
that names a problem, translations and wording improvements become test
changes. The customer workspace begins with identities because identities are
the anchors that let presentation change safely.

\begin{lstlisting}[caption={The committed edit value.}]
@SwingForm(id = "customer/edit")
public record CustomerEdit(
        @TextField
        @Required
        @Length(max = 80)
        String name,

        @TextField
        @DecimalRange(min = "0.00")
        BigDecimal creditLimit,

        @CheckBox
        boolean active) {
}
\end{lstlisting}

The record is not the live Swing form. It is the committed value shape. It is
the object that can be saved, compared, audited, and passed to domain code
after editing succeeds. The generated form model will own editing state: raw
text, parsed values, parse errors, validation problems, dirty state, touched
state, and guarded result creation. This distinction keeps the record small
and immutable while giving the UI a proper place for unfinished input.

Records are pleasant enough that developers may be tempted to use them as live
state. Resist that in an editor. A partially typed decimal is not a
\api{BigDecimal}. A text field that has been touched but not changed is not a
new committed value. A validation problem is not a property of the record; it
is a property of the current edit attempt. The form model is the transaction
boundary between widgets and records.

\begin{lstlisting}[caption={The table row value.}]
@SwingTable(id = "customer/table")
public record CustomerRow(
        @Column(width = 80)
        CustomerId id,

        @Column(width = 220)
        String displayName,

        @Column(width = 100)
        CustomerStatus status,

        @Column(width = 120)
        LocalDate lastOrderDate) {
}
\end{lstlisting}

Define the row as a value before constructing the table. A row record should
contain the values the table displays and the identity needed for selection
and refresh. Generated table companions can then expose table keys, column
keys, descriptors, resource keys, and Swing model helpers. The \api{JTable}
comes later, after the row contract exists.

This order prevents the table from becoming the schema. A \api{JTable} has
view indexes, model indexes, sorters, renderers, localized headers, column
models, and selection models. Those details are necessary at runtime, but they
are poor places to store meaning. The generated or handwritten descriptor
names the row type and columns before Swing displays them. When the user moves
a column, the descriptor identity does not move. When a locale changes the
header, the persistence id does not change. When a test asserts selection by
row value, it need not know the current view row.

\begin{lstlisting}[caption={Actions are declared before buttons.}]
final class CustomerWorkspacePresenter {
    @UiAction(
            id = "customer/refresh",
            text = "customer/refresh/text",
            tooltip = "customer/refresh/tooltip")
    void refresh() {
        loadCustomers();
    }

    @UiAction(
            id = "customer/save",
            text = "customer/save/text",
            tooltip = "customer/save/tooltip",
            acceleratorKey = KeyEvent.VK_S,
            acceleratorModifiers = InputEvent.CTRL_DOWN_MASK)
    void save() {
        saveCurrentCustomer();
    }
}
\end{lstlisting}

Commands should be declared before buttons, menu items, and shortcuts. The
workspace has operations such as Refresh, Save, Reset, Delete, and Help. Each
operation has one command identity and one enabled-state policy. Generated
action companions can create descriptors and command factories. The view
adapts commands to buttons and key bindings. It does not own Save logic.

The same identity rule applies to resources and help. A label resource key
names visible text for a field. A help topic names a documentation target. A
component key names a view contract for tests and focus. These keys should be
near generated companions or handwritten descriptors, not sprinkled through
component construction. A generated form descriptor can carry label, tooltip,
problem, and help metadata for a field. A generated table descriptor can carry
column headers, tooltips, help topics, widths, and persistence ids. A generated
action descriptor can carry command text, tooltip, mnemonic, accelerator, and
selectable state.

\begin{principlebox}{Facts First}
Declare stable facts before constructing visible components. A screen that
starts with keys, descriptors, models, and commands can change layout and text
without losing meaning.
\end{principlebox}

This does not mean every screen begins with an annotation. A one-off custom
visual editor may declare handwritten keys and descriptors. A mature legacy
table may be migrated one column at a time. The rule is broader than
generation: do not let Swing component instances be the only source of truth
for application facts. Corusco generation is a convenient way to make that
rule cheap for common forms, tables, and actions.

The stable facts should also be grouped by owner. Form field keys belong with
the form source or form companion. Table column keys belong with the table row
source or table companion. Action keys belong with the presenter or action
companion. Component keys may belong with the view because they name view
contracts rather than domain facts. Resource keys may be generated from the
same source declarations but resolved by the application resource service.
Help topics may be generated for fields and columns or declared in a feature
help companion. A key used everywhere but owned nowhere becomes difficult to
rename and difficult to review.

A useful early design exercise is to write the screen's vocabulary before the
screen's component tree. Which rows exist? Which columns are durable? Which
fields are editable? Which operations can the user perform? Which help topics
are meaningful? Which resources must be localized? Which facts are persisted
between sessions? This vocabulary need not be perfect. It will change during
implementation. But it gives the screen a semantic skeleton before Swing
components, listeners, and layout constraints multiply.

\section{Presenter state and command vocabulary}

The presenter owns screen state that is not purely visual: observable row
source, selected row or selected customer id, form model for the detail side,
command set, busy state, task service, status text, problem sets, and lifecycle
scopes. This does not mean the presenter should become a god object. It means
state needs a non-widget owner. Split services and models when they become
concepts of their own.

The presenter is the place where screen facts become workflow. It knows that a
selected row should load detail, that Save depends on dirty state and
validation, that Refresh starts a background task, and that a closed view
should cancel view-owned work. It should not know how MigLayout aligns a
label, and it should not paint a table cell. If a service call, permission
rule, or validation rule becomes large, move it into a service or model. The
presenter coordinates; it should not absorb the whole application.

\begin{lstlisting}[caption={Presenter state sketch.}]
final class CustomerWorkspacePresenter implements AutoCloseable {
    private final ObservableArrayList<CustomerRow> rows =
            ObservableArrayList.empty();
    private final SimpleValue<CustomerRow> selectedRow =
            SimpleValue.empty();
    private final SimpleValue<String> statusText =
            SimpleValue.of("Ready");
    private final SimpleValue<Boolean> refreshBusy =
            SimpleValue.of(false);
    private final TaskService tasks =
            SwingTaskServices.virtualThreads();
    private final SubscriptionScope scope =
            new SubscriptionScope();

    ObservableList<CustomerRow> rows() {
        return rows;
    }
}
\end{lstlisting}

The class sketch is incomplete by design. The important part is ownership. If
rows change, inspect the row source. If selection behaves oddly, inspect the
selected-row value and the selection binding. If Save is enabled at the wrong
time, inspect the derived command state. If a background result arrives too
late, inspect the generation counter and task callback. Without these named
owners, the reviewer must search through listeners and component state hoping
to find the one place where a fact is currently stored.

Commands are part of presenter state because they express operations in terms
of screen state. The Save command is enabled when the form is dirty, valid,
editable, and not saving. The Delete command is enabled when a durable
selection exists and permission allows deletion. The Refresh command may stay
enabled even while the form is invalid because it does not depend on form
validity. These are not button rules. They are operation rules.

\begin{lstlisting}[caption={Command state derived from screen facts.}]
DerivedValue<Boolean> canSave = DerivedValue.of(
        () -> formDirty.value()
                && !formProblems.value().hasErrors()
                && !saveBusy.value(),
        formDirty, formProblems, saveBusy);

scope.add(canSave.subscribe(event ->
        saveCommand.setEnabled(Boolean.TRUE.equals(event.newValue()),
                event.origin())));
\end{lstlisting}

The generated action descriptor supplies identity and presentation metadata.
The presenter supplies runtime availability. This distinction lets a toolbar
button, menu item, popup item, key binding, and command palette observe the
same command. If the keyboard shortcut works while the Save button is disabled,
there are probably two operation paths. The fix is not another boolean; the
fix is one command.

Status text deserves a model too. A status label that is set directly by every
listener becomes a race between concerns. Refresh may set "Loading". Save may
set "Saved". Validation may set "Fix validation errors". A background failure
may set "Network unavailable". Decide whether the screen has one status value,
a prioritized status model, or several status regions. Then bind labels to
that state. Do not leave status policy as a collection of calls to
\api{setText}.

The same applies to disabled reasons. Save disabled because the form is clean
is different from Save disabled because validation failed, permission is
missing, or a save is already running. A boolean command enabled state is
necessary for Swing, but it may not be sufficient for users. A neighboring
reason value can feed a tooltip, status line, command palette, or accessible
description. The presenter is the place to derive that reason from form,
permission, and task state.

Presenter state should be testable without Swing where possible. A core test
can construct the presenter with fake services, load deterministic rows,
select a row, edit the form model, and assert command state. EDT tests should
cover bindings and component behavior. This separation keeps the expensive
Swing tests focused and lets workflow rules run quickly.

Permissions belong in this layer too. A Delete command may be disabled because
there is no selection, because the selected customer is archived, because a
delete task is already running, or because the current user lacks permission.
Those reasons have different user-facing consequences. A permission failure
may need a help link or audit message. A missing selection may need no
explanation. A busy task may need a progress indicator. Modeling only one
boolean forces the view to invent text it cannot justify. Presenter state can
derive both enabled booleans and explanatory values from the same facts.

Undo and reset are a useful test of command vocabulary. Reset detail form
should restore the current form model to the loaded customer. Undo may be a
broader application command that reverses the last edit. Cancel may close a
dialog or discard workspace changes depending on container. These operations
sound related, but they are not the same command. Naming them precisely in the
presenter prevents the view from turning every "go back" gesture into a local
button listener.

\section{View assembly with MigLayout}

The view is an ordinary Swing panel. Use MigLayout for the main geometry. Use
FlatLaf for modern styling. Use a split between the table/search area and the
detail form. Keep layout visible and reviewable. Do not bury the main screen
in decorative cards. This is an operational desktop screen. Density and
predictable scanning matter more than a marketing layout.

\begin{lstlisting}[caption={Workspace layout shape.}]
JPanel panel = new JPanel(new MigLayout(
        "fill, insets 18, gap 12",
        "[grow, fill][320!, fill]",
        "[][grow, fill][]"));

panel.add(searchBar, "span 2, growx, wrap");
panel.add(new JScrollPane(customerTable), "grow");
panel.add(detailPanel, "growy, wrap");
panel.add(statusLabel, "span 2, growx");
\end{lstlisting}

The layout should support repeated work. The search area belongs where the
user begins. The table should receive enough space to compare rows. The detail
panel should keep labels and fields aligned. The status line should be visible
without competing with the form. MigLayout lets the source describe these
priorities directly: which region grows, which width is fixed, where wrapping
occurs, and which gaps express related or unrelated content.

The detail panel implements the generated view contract. The application still
constructs fields and arranges them. The generated behavior plan installs text
bindings, selected-state bindings, validation tooltips, validation borders,
help topics, and accessibility metadata. The view contract is a narrow bridge,
not a replacement for the panel.

\begin{lstlisting}[caption={A handwritten panel implements a generated view contract.}]
final class CustomerDetailPanel extends JPanel
        implements CustomerEditView {
    private final JTextField name = new JTextField();
    private final JTextField creditLimit = new JTextField();
    private final JCheckBox active = new JCheckBox();

    CustomerDetailPanel() {
        super(new MigLayout(
                "fillx, insets 12, gap 10",
                "[][grow, fill]",
                "[][][]push[]"));
        add(new JLabel("Name"));
        add(name, "wrap");
        add(new JLabel("Credit limit"));
        add(creditLimit, "wrap");
        add(new JLabel("Active"));
        add(active, "wrap");
    }

    @Override
    public JTextField nameField() {
        return name;
    }
}
\end{lstlisting}

The source listing is intentionally mundane. The panel is allowed to be
ordinary. It arranges labels and editors in a way a Swing developer can
inspect. It implements the generated contract so generated bindings can find
the components they need. If the design changes, change the MigLayout source.
If a generated binding changes, inspect the generated behavior plan. The two
concerns meet at the contract methods.

Generated plans should be read, not merely trusted. They are ordinary Java
source produced at compile time. A developer can open the generated view
contract, see which components are expected, and inspect which bindings or
behaviors are installed. This inspectability is one reason compile-time
generation fits a Swing codebase better than a hidden runtime scanner. When
the screen misbehaves, the generated source becomes documentation for the
agreement between annotated model and handwritten view.

\begin{lstlisting}[caption={Generated plan installed into a handwritten view.}]
BehaviorScope behaviors = new BehaviorScope(helpService);

CustomerEditBehaviorPlan.install(
        detailPanel,
        presenter.formModel(),
        behaviors);
\end{lstlisting}

The exact generated facade names depend on the source type. The pattern is the
same: handwritten view, generated synchronization plan, lifecycle-owned scope.
Installing bindings creates listener relationships. Closing the scope must
detach them. A real screen will be opened, closed, replaced, and perhaps
embedded in different containers. Binding lifecycle is not optional.

\begin{figure}[htbp]
\centering
\begin{minipage}[t]{0.48\linewidth}
\includegraphics[width=\linewidth]{\bindingFormFigure}
\end{minipage}\hfill
\begin{minipage}[t]{0.48\linewidth}
\includegraphics[width=\linewidth]{\behaviorPlanFigure}
\end{minipage}
\caption{The handwritten view owns geometry, while generated form metadata and behavior plans own routine synchronization, validation presentation, help, and cleanup.}
\label{fig:bookapp-form-behaviors}
\end{figure}

Figure~\ref{fig:bookapp-form-behaviors} is the key compromise. Swing
components remain visible in the code. MigLayout remains the geometry
language. Corusco supplies field models, descriptors, and binding plans.
Neither side pretends to be the other. This is why the screen can be both
generated enough to avoid repetition and handwritten enough to remain designed.

FlatLaf belongs at the application boundary. The example code may call a
support helper that installs a FlatLaf theme before screenshots. Production
applications usually set the Look-and-Feel at startup. Individual panels
should not repeatedly install the LAF. They should use standard Swing
components, client properties only where needed, and layout constraints that
work under light and dark themes. Screenshots should verify both density and
readability.

Focus order should be reviewed with the same seriousness as layout. A user
will tab through the search field, table, detail fields, and command row. The
order should follow the workflow, not construction accidents. Generated view
contracts and component keys help tests find components; they do not decide
focus traversal by themselves. Swing's default focus traversal is often good
when components are added in logical order. When the visual layout and add
order diverge, review focus explicitly.

Resize behavior is part of the design. The table should normally absorb most
extra horizontal and vertical space. The detail panel may have a fixed
comfortable width. The status line should grow horizontally but not vertically.
The command row should keep natural button sizes. A screenshot at one size is
not enough. The layout should be inspected at a narrower width, a taller
workspace, and a high-DPI scale. MigLayout makes these policies visible in
constraints, but it cannot choose them for the application.

\section{Tables and master-detail selection}

The customer table uses the generated descriptor and an observable row source.
The table model is installed on the EDT. Selection binding writes selected row
state. A table-state controller restores column layout and sort state. This is
more code than \api{new JTable(data, headers)}, but it has stable identity,
selection policy, table persistence, and lifecycle cleanup.

\begin{lstlisting}[caption={Table assembly shape.}]
EdtObservableList<CustomerRow> swingRows =
        EdtObservableList.of(presenter.rows());

ObservableTableModel<CustomerRow> model =
        ObservableTableModel.of(swingRows,
                CustomerRowTableDescriptor.DESCRIPTOR);

customerTable.setModel(model);

scope.add(TableSelectionBinding.bind(
        customerTable, model,
        presenter.selectedModelRow(),
        presenter.selectedRow()));

scope.add(TableStateController.install(
        customerTable, model, tableStateStore));
\end{lstlisting}

The extra code is where table meaning becomes explicit. The row list is
observable. The table descriptor supplies columns. The selection binding states
which presenter values receive selected coordinates and row values. The
table-state controller states where column layout is restored and saved. In a
plain \api{new JTable(data, headers)} table, these facts still exist; they are
merely implicit, often as indexes and listeners that will be harder to find
later.

\begin{figure}[htbp]
\centering
\begin{minipage}[t]{0.48\linewidth}
\includegraphics[width=\linewidth]{\coruscoTableDescriptorFigure}
\end{minipage}\hfill
\begin{minipage}[t]{0.48\linewidth}
\includegraphics[width=\linewidth]{\coruscoTableStateFigure}
\end{minipage}
\caption{The table descriptor names columns before Swing displays them; table state then persists user layout against stable ids rather than localized headers or view indexes.}
\label{fig:bookapp-table-state}
\end{figure}

Figure~\ref{fig:bookapp-table-state} explains why descriptor-backed tables
matter in a complete screen. The row record and descriptor define meaning. The
\api{JTable} displays it. The column model lets the user rearrange it. The
state controller persists the arrangement. If meaning lives only in column
indexes, these steps become tangled. If meaning lives in column keys and
descriptors, view order can change without changing identity.

When the selected row changes, the detail form should load or display the
corresponding customer. This is a master-detail relationship. Avoid putting
the whole workflow in a table selection listener. Selection binding writes
selected state. Presenter logic reacts to selected state. Detail loading uses
tasks or master-detail values as appropriate.

\begin{lstlisting}[caption={Selection reaction in presenter state.}]
scope.add(selectedRow.subscribe(event -> {
    CustomerRow row = event.newValue();
    if (row == null) {
        clearDetail();
    } else {
        loadDetail(row.id());
    }
}));
\end{lstlisting}

A production screen usually has at least two selection speeds. Selecting a row
should feel immediate: the visible selection, command state, and perhaps a
small amount of cached detail can update at once. Loading expensive detail may
take longer and must be treated as asynchronous. The selected row value is
therefore not the same as the loaded detail value. Keeping them separate lets
the UI show "loading" for the current selection while preventing an older load
from overwriting the current detail panel.

Selection preservation during refresh deserves policy. If the refreshed row
set still contains the selected customer id, the screen may preserve
selection. If the selected customer disappeared, the screen may clear detail,
show a message, or select the nearest row. Do not let the table sorter decide
this accidentally. The presenter should compare durable row ids, not view
indexes. Generated row identities and column descriptors make that comparison
straightforward.

Filtering and sorting add another coordinate boundary. A selected view row is
not a selected model row. A selected model row is not necessarily a selected
domain id. A table descriptor does not remove this distinction; it makes it
safer to cross. Use table selection bindings that explicitly convert
coordinates and write selected state. Avoid reading \api{getSelectedRow} in
every command handler.

Renderers remain explicit Swing code. A generated descriptor can supply value
type and column identity, but a status badge, risk color, date format, or
validation marker needs renderer policy. Keep renderers small, stateless, and
aware of performance. Let the descriptor identify the column; let the renderer
present the cell. Do not put presenter state into renderer components.

Large-data behavior should be considered even when the example data is small.
The customer table may eventually show tens of thousands of rows or a virtual
slice of a larger result. The descriptor-backed model helps avoid column
switch statements, but row loading, filtering, sorting, and update precision
still matter. Use observable list updates that describe what changed. Avoid
replacing the entire table model for a small row edit. Preserve selection by
durable id where possible. Repaint the rows that changed, not the whole table,
unless the whole table truly changed.

Export and copy actions should use the same column identities. If the user
exports selected customers, the export should not scrape visible table headers
and view indexes as its only schema. The table descriptor can name columns,
value extractors, resource keys, and visibility policy. The export command can
decide whether to use all columns, visible columns, or a separate export
descriptor. This is another example of stable facts becoming useful after the
first visible table works.

\section{Forms, validation, resources, and help}

The detail form edits the selected customer. Its generated form model owns raw
text, parsed values, validation, dirty state, and committed result creation.
The view owns components and layout. The presenter owns when a form is loaded,
saved, reset, or discarded. This separation prevents the common Swing mistake
where the current text field values are both the user interface and the
application state.

Validation problems from the form model appear in field borders, tooltips,
summary labels, command enablement, and tests. The same \api{ProblemSet}
supports all of these. Do not recalculate validation separately for each
surface. A field with invalid input should not tell the tooltip one thing, the
Save command another thing, and the problem summary a third thing.

\begin{lstlisting}[caption={Save enablement from form state.}]
DerivedValue<Boolean> canSave = DerivedValue.of(
        () -> formDirty.value()
                && !formProblems.value().hasErrors()
                && !saveBusy.value(),
        formDirty, formProblems, saveBusy);
\end{lstlisting}

The generated or handwritten command state observes this value. The button,
menu item, shortcut, and command palette observe the command. This indirection
is not ceremony. It is how a screen prevents a keyboard shortcut from saving
invalid data while the visible Save button is disabled.

Validation presentation should be staged. A newly loaded form may contain
empty optional fields without screaming at the user. A field being edited may
show parse feedback before full validation. A Save attempt may reveal all
blocking errors. The form model supplies problems and field identities; the
presentation policy decides when to show them. Do not confuse the existence of
a problem with the decision to draw it loudly.

Resources and help should be generated or keyed before they are painted.
Button text, labels, tooltips, table headers, validation messages, and help
indicators come from resources and generated keys. F1 help uses generated help
topics or handwritten topics attached to descriptors. The screen should not
contain visible strings as stable ids. This matters for the book as well:
screenshots should remain visually stable, but source should still demonstrate
resource separation.

Resource separation also keeps examples from teaching bad habits accidentally.
It is acceptable for a compact book example to use a small in-memory resource
map, but components should still be driven by resource keys where the real
application would be. This gives the reader the habit of treating visible text
as presentation, not identity. The same principle applies to help topics and
tooltips. A tooltip may contain static help, validation feedback, and disabled
reason text; a structured tooltip policy is easier to test and compose than a
sequence of calls that overwrite \api{setToolTipText}.

Help routing is part of the screen, not an afterthought. F1 on the name field
should open help for the name field, not a generic application page. F1 on a
table column should open column help if the descriptor supplies it. F1 on a
disabled Save button may open help for the Save operation and explain the
disabled reason. Generated help topics and descriptors make this possible
without a brittle map from English text to help pages.

Accessibility follows the same metadata. A label resource can contribute to an
accessible name. A tooltip or help resource can contribute to an accessible
description. A disabled reason can explain why a command is unavailable. A
component key can help tests and focus resolvers find the correct component.
Do not bolt accessibility onto the finished screen by copying visible text into
random properties. Let the same descriptors and resources feed visual and
accessible presentation where appropriate.

Save and cancel workflows deserve special attention. A form edit should commit
the active editor before validation is judged. If the user is typing in a text
field and presses Ctrl+S, the document value and form model should agree.
Generated text bindings help, but custom editors and formatted fields may need
explicit commit behavior. Cancel should distinguish a clean form from a dirty
form. Dirty cancel confirmation belongs to workflow policy, not to the text
field that happened to change first.

Problem summaries are optional but often useful. A dense detail form may show
field-level markers and tooltips, but a summary gives keyboard and screen
reader users a stable place to review all blocking issues. The summary should
read from the same \api{ProblemSet} as field decoration and command
enablement. It should not run a second validation pass with different wording.
If selecting a problem in the summary focuses the field, use field keys and
component keys to route focus rather than comparing label text.

\section{Background work and busy state}

Refresh is a command that starts background work. The task body loads rows.
The success callback applies rows to the observable list. Busy state disables
incompatible commands and drives an overlay or status text. The task handle is
owned by a lifecycle scope so closing the screen can cancel work or detach
callbacks.

\begin{lstlisting}[caption={Refresh as task-owned workflow.}]
void loadCustomers() {
    GenerationCounter.Generation generation =
            refreshGenerations.advance();
    TaskHandle<List<CustomerRow>> handle = tasks.submit(
            cancellation -> customerService.loadRows(cancellation),
            TaskCallbacks.onSuccess(loaded ->
                    refreshGenerations.tryAccept(generation, loaded, rows -> {
                        this.rows.batch(list -> {
                            list.clear();
                            rows.forEach(list::add);
                        });
                        statusText.setValue(
                                "Loaded " + rows.size() + " customers");
                    })));
    scope.add(handle);
}
\end{lstlisting}

The example assumes callbacks are delivered on the EDT through
\api{SwingTaskServices}. If a different service is used, the EDT boundary must
be explicit. A background thread may load rows. It must not mutate Swing
components directly. A callback that updates an observable list adapted to
Swing should do so on the intended thread. The lifecycle chapter explains the
mechanics; the screen-building rule is to make the boundary visible.

Refresh has more policy than its name suggests. Duplicate refreshes may be
blocked, queued, or allowed to cancel previous work. Old rows may remain
visible while new rows load, or the table may clear immediately. Failures may
show a status message, problem banner, retry command, or modal dialog.
Selection may be preserved by customer id, cleared, or moved to the nearest
row. Corusco does not answer these business questions. It provides values,
lists, task handles, generation guards, and scopes so the application can
answer them in readable code.

Busy state affects more than one component. The Refresh command disables. The
Save command may disable during save. The status line changes. The table area
may show a busy overlay. Because busy is modeled, these presentations are
bindings, not duplicate booleans.

\begin{lstlisting}[caption={Busy overlay in the workspace view.}]
JLayer<JComponent> tableLayer = new JLayer<>(tablePanel);
scope.add(BusyOverlayBinding.install(
        tableLayer, presenter.refreshBusy()));
\end{lstlisting}

Use the overlay when input should be blocked. Do not use it as decoration when
the user should continue interacting. A busy veil that visually blocks a table
while the table still accepts selection is misleading. A busy overlay that
steals focus from a field during save may be worse than no overlay. Busy
presentation is workflow policy, not just painting.

\begin{figure}[htbp]
\centering
\begin{minipage}[t]{0.48\linewidth}
\includegraphics[width=\linewidth]{\busyTaskFigure}
\end{minipage}\hfill
\begin{minipage}[t]{0.48\linewidth}
\includegraphics[width=\linewidth]{\backgroundStatusFigure}
\end{minipage}
\caption{Background work must surface through command state, busy overlays, status/progress values, cancellation, and stale-result suppression rather than direct component mutation from worker code.}
\label{fig:bookapp-background}
\end{figure}

Figure~\ref{fig:bookapp-background} shows why background work belongs in the
screen architecture, not in one button listener. A task affects commands,
status, progress, cancellation, and stale-result policy. If those facts are
modeled, the UI can present them consistently. If they live in a worker's
anonymous callback, every new busy presentation will copy a little more state.

Virtual threads are useful for blocking service calls, but they do not remove
the Swing rule. The task may run on a virtual thread. The Swing mutation must
return to the EDT. Structured concurrency may help organize service-side work.
It does not make a \api{JTable} thread-safe. The modern Java chapter explains
the boundary; the workspace applies it.

Failure handling should be modeled before it is painted. A refresh failure may
leave old rows visible with a stale warning. A save failure may keep the form
dirty and show a retryable error. A permission failure may disable a command
and route to help. A validation failure is not a task failure at all. If all
failures are caught by one callback that shows a modal error dialog, the user
will receive poor guidance and tests will have little to assert. Treat failure
as data first, then decide how the view presents it.

Cancellation should have a user story. Closing the workspace should cancel
view-owned loads or at least detach callbacks. Pressing Refresh again may
cancel the previous refresh or let it finish but suppress its stale result.
Saving may be noncancelable after the server accepted the request. These
policies differ. A task handle and cancellation token let the screen express
them. A raw worker thread hidden in a listener does not.

\section{Lifecycle and cleanup}

A real screen closes. It may also be hidden, re-opened, embedded in another
shell, or replaced after navigation. Every listener, binding, table-state
controller, task handle, generated behavior plan, and adapter needs an owner
that closes with the screen. Corusco gives several lifecycle handles:
\api{Disposable}, \api{SubscriptionScope}, \api{DetachableScope},
\api{BehaviorScope}, and task handles. Use them intentionally.

There are at least three lifecycles in the workspace. The application
lifecycle owns services and long-lived stores. The workspace lifecycle owns
the presenter, row source, command set, and screen-level tasks. The view
lifecycle owns Swing bindings, table controllers, generated behavior plans,
and adapters installed into components. Mixing these lifecycles is a common
leak source. A command may outlive a dialog; a Swing action adapter installed
in that dialog should not.

\begin{lstlisting}[caption={View lifecycle owns Swing installations.}]
final class CustomerWorkspaceView implements Disposable {
    private final BehaviorScope behaviors = new BehaviorScope();
    private final BindingScope bindings = new BindingScope();

    void install(CustomerWorkspacePresenter presenter) {
        CustomerEditBehaviorPlan.install(
                detailPanel, presenter.formModel(), behaviors);
        bindings.add(TableStateController.install(
                table, presenter.tableModel(), tableStateStore));
    }

    @Override
    public void close() {
        bindings.close();
        behaviors.close();
    }
}
\end{lstlisting}

The code shape is more important than the exact class names. The view that
installs Swing relationships closes them. The presenter that starts
presenter-owned subscriptions closes them. The task service that owns
background execution has application or screen lifetime according to policy.
No listener should have to guess who will remove it.

Lifecycle also affects generated code. Generated behavior plans install
bindings and behaviors into a provided scope. They should not hide global
singletons or static listener registries. Generated command adapters should be
closed with the view surface that installed them. Generated table models
adapted to Swing should be detached when the table no longer needs them. A
generated companion that makes cleanup harder is not removing wiring; it is
relocating it.

Closing policy should be tested. Create a view, install bindings, close the
scope, change the model, and assert that disposed components no longer update.
Start a task, close the presenter, and assert that callbacks do not mutate the
closed view. Install a key binding, close the behavior scope, and assert that
the previous input-map entry is restored. These are not glamorous tests, but
they prevent the slow accumulation of leaks and stale callbacks.

\section{Testing and screenshots}

Use layered tests. Core tests cover form models, validation, commands, values,
and lists. Generated-source tests cover annotations and companions. EDT tests
cover bindings, table selection, and table state. Task tests cover cancellation
and stale results. Screenshot smoke tests cover the companion visual shell.
One large end-to-end test cannot replace these layers because it cannot tell
the reader which boundary failed.

\begin{lstlisting}[caption={Smoke-test the companion window.}]
SwingUtilities.invokeAndWait(() -> {
    JInternalFrame frame = BookWorkspaceExample.createWindow();
    assertThat(frame.getTitle()).isEqualTo("Book workspace");
    assertThat(frame.getContentPane().getComponentCount())
            .isGreaterThan(0);
});
\end{lstlisting}

The current companion test exercises construction. Deeper tests should target
the richer generated examples and core models. As the screen grows, tests
should grow sideways rather than only upward. Add a core test for a validation
rule. Add a command test for enablement. Add an EDT test for a binding. Add a
table-state test for persisted columns. Add a task test for cancellation and
stale results. Keep one or two screenshot smoke tests as visual anchors. This
gives a reader, reviewer, and future maintainer several small proofs instead
of one large performance.

\begin{figure}[htbp]
\centering
\begin{minipage}[t]{0.48\linewidth}
\includegraphics[width=\linewidth]{\testingKeysFigure}
\end{minipage}\hfill
\begin{minipage}[t]{0.48\linewidth}
\includegraphics[width=\linewidth]{\testingHarnessFigure}
\end{minipage}
\caption{A complete screen should be tested through stable component and command identities, while screenshots verify the rendered FlatLaf/MigLayout result.}
\label{fig:bookapp-testing}
\end{figure}

Figure~\ref{fig:bookapp-testing} shows the testing compromise. Stable keys
and generated contracts make tests robust. Screenshots make visual regressions
visible. The harness should locate components through keys or contracts, not
through the text whose visual rendering it is trying to inspect. If a label is
copyedited, the screenshot diff should show the text change; the test harness
should not fail to find the field.

Screenshot source should be deterministic: fixed size, FlatLaf, seeded data,
stable content, and a known state. The compact \bookAppExample{} exists partly
for this reason. It gives the book a visual anchor without forcing every
architectural feature into one unreadable listing. As the example grows, keep
screenshot code and model tests aligned. A screenshot program that is not
exercised by any test drifts quickly.

The screenshot should show the reader what the prose is discussing. It should
not be a glamour image detached from the actual example. The table should be
populated by the same deterministic data used by tests where practical. The
selected row, validation state, and busy indication should be chosen because
they teach a concept in the surrounding text. If the screenshot requires many
special setup calls that no application would use, the example is probably
serving the figure more than the reader.

For book work, screenshot acceptance is part of documentation acceptance. A
chapter that discusses visual behavior should include enough figures to let
the reader see the state. The real-screen chapter includes the normal
workspace, busy-detail state, generated form behavior, table state,
background-progress state, and testing harness context. The figures should be
generated from companion examples or shared figure harnesses so they can be
refreshed when the code changes.

\section{Construction order}

A practical construction order keeps the system from starting as a component
tree full of hidden state. First define records and stable ids. Then add
annotations for form, table, actions, validation, and help. Build core
presenter state. Generate and inspect companions. Create the MigLayout view.
Install table model, selection binding, and table state. Install form bindings
and behaviors. Add commands and keyboard shortcuts. Add task service, busy
state, overlays, and stale-result suppression. Add layered tests and
screenshot harness output.

This order is forgiving. After records and stable ids exist, generation can be
inspected before the view is finished. After presenter state exists, command
and validation tests can be written before the final layout. After the
MigLayout view exists, screenshots can guide visual refinement while core
tests continue to protect behavior. Building in this order gives the developer
frequent checkpoints and prevents the late discovery that important facts were
embedded only in component positions or localized strings.

The order also helps migration. A legacy screen may already have a working
view. Do not throw it away. Introduce stable row and field identities. Replace
index-based table code with descriptors. Introduce commands behind existing
buttons. Move validation into a form model. Add table state by stable ids.
Wrap background work in task services. Add lifecycle scopes around listeners.
Each step should reduce one class of hidden coupling while keeping the screen
recognizable to users.

\begin{migrationbox}{Screen Migration}
Migrate by ownership, not by visual region. A row descriptor, command model,
form model, or task boundary can improve a legacy screen without redesigning
the whole view at once.
\end{migrationbox}

Some work should remain handwritten from the beginning: layout, service calls,
navigation, permission checks, custom renderers, business workflows, unusual
validation rules, error presentation policy, and domain-specific status text.
Corusco removes repeated wiring. It does not remove design. That sentence is
the best defense against both overuse and underuse. Use the library where
repetition hides a stable concept. Keep application code where the domain
decision is genuinely specific.

\begin{examplebox}{Book Application}
\bookAppExample{} assembles the book's preferred screen shell with MigLayout,
FlatLaf-oriented construction, deterministic table rows, and a detail form. It
is the screenshot-oriented seed for the larger customer workspace described in
this chapter.
\end{examplebox}

\section{A user story through the architecture}

Follow one ordinary user story through the screen. The workspace opens. The
view constructs components on the EDT, sets up MigLayout, creates the table,
creates the detail panel, and installs generated form behavior into a
lifecycle scope. The presenter constructs row state, selected-row state,
commands, task handles, and status values. The table model is installed from
the generated descriptor. Table state is restored by stable ids. No service
call has happened yet, but the screen already has owners for layout, rows,
selection, commands, bindings, and cleanup.

The user presses Refresh. The toolbar button, menu item, and shortcut all
invoke the same Refresh command. The command calls the presenter method. The
presenter advances a generation counter, sets refresh-busy state, starts a
task, and records the task handle in a scope. The view observes busy state and
shows the busy presentation selected for this screen. The Refresh command may
disable or remain enabled according to policy. The status value changes to
"Loading customers" through a model value, not by reaching into a label from
the task body.

The service call runs away from the EDT. It may use a virtual thread, a pooled
thread, or a test fake. It receives a cancellation token. It does not mutate
Swing components. If the user closes the workspace while the call is running,
the workspace lifecycle closes the task handle or at least detaches the
callback. If the user presses Refresh again, policy decides whether the older
task is cancelled or merely made stale. The generation counter makes that
policy explicit.

Rows arrive. The success callback returns to the EDT boundary promised by the
task service. The presenter checks whether the generation is still current. If
it is stale, the result is ignored. If it is current, the row list is updated
with precise observable-list changes. The table model observes the row list
and fires table events. The table repaint is a consequence of model update,
not a direct call from the worker. The status value changes to "Loaded 125
customers" or equivalent text derived from the result.

The user selects a row. Swing selection changes first, but selection binding
translates the view coordinate to model coordinate and writes selected state
into presenter values. Commands that depend on selection update. The detail
load workflow sees the selected customer id and starts a detail load or uses
cached detail. If the table is sorted, the selection binding crosses the
coordinate boundary in one place. Command handlers do not ask the table for
the selected view row.

The detail arrives. The presenter creates or replaces the generated form model
for the selected customer. The view binding scope connects the generated form
model to the existing detail components. The name field shows committed text.
The credit limit field shows formatted raw text. The active check box shows
the boolean value. The form is clean. Save is disabled because there are no
changes. The disabled reason may be empty or "No changes" depending on product
policy.

The user edits the credit limit and types an invalid decimal. The document
binding updates the text-field model. The converter records parse failure.
The form problem set changes. The field decoration updates through validation
behavior. The tooltip policy composes the parse or validation message with
static help. The Save command remains disabled because the form has errors.
The problem summary, if present, reads the same problem set. No component
needs to run its own validation rule.

The user fixes the value. Parse state becomes valid. The form remains dirty.
The problem set clears. The Save command becomes enabled because dirty,
valid, editable, and not busy are all true. If the command is visible in a
button, menu item, and keyboard shortcut, all surfaces receive the same state.
The view did not enable the button directly. The command did.

The user presses Ctrl+S. The key binding invokes the Save command. Before the
command judges validity, active-editor commit policy ensures that the current
editor value has reached the model. The presenter starts a save task or calls
a service according to policy. Save busy state disables incompatible commands.
The status line changes. If saving is asynchronous, stale-result policy
protects against a late response for a customer that is no longer selected.
If saving succeeds, the committed record becomes the new dirty baseline and
the command disables because the form is clean. If saving fails, failure state
is recorded and presented without losing the user's edits.

The user asks for help. F1 on a field uses generated field help metadata. F1
on a table column uses generated column help metadata plus current column
context. F1 on Save uses command help or action identity. The help service
receives a structured request. It may open an embedded help panel, browser, or
test recorder. It does not infer the help topic from visible English text.

The user closes the workspace. The view scope closes generated behavior
bindings, table controllers, key bindings, tooltips, adapters, and selection
bindings. The presenter scope closes subscriptions and task handles it owns.
Long-lived services remain alive if they belong to the application. The
screen's closure is therefore not a best-effort disposal hidden in a window
listener. It is the inverse of installation.

This story may sound verbose, but each step is small in code when the owners
exist. It becomes verbose only when the architecture is absent. Then Refresh
has a worker, selection has a listener, detail has another listener, Save has
a button callback, Ctrl+S has a root-pane action, validation has document
listeners, help has component names, and close has a window listener that
knows none of the above. The purpose of Corusco in a real screen is to keep
the ordinary story ordinary.

\section{File and package shape}

A maintainable screen should be easy to browse in the project tree. The exact
package names depend on the application, but the file roles should be visible.
A small feature package might contain \api{CustomerRow}, \api{CustomerEdit},
\api{CustomerWorkspacePresenter}, \api{CustomerWorkspacePanel},
\api{CustomerWorkspaceResources}, and \api{CustomerWorkspaceTest}. Generated
companions appear beside annotated sources or under the generated-source root.
Screenshot fixtures live with book or visual-test support. Avoid a package
named \api{ui} that contains every panel, renderer, presenter, resource map,
and test helper for the whole application.

The model-facing files should be boring. \api{CustomerRow} describes table
rows. \api{CustomerEdit} describes committed form values.
\api{CustomerStatus} describes status choices. If the screen has a durable
customer id, it should have a small type or at least a clearly named value.
These files should not import Swing. They may import Corusco annotations if
they are generation inputs. They may import validation annotations if the
constraints are part of the UI contract. They should not know about buttons,
tables, or layout.

The presenter-facing files should read like workflow. A presenter constructor
receives services, resources, task services, or stores it needs. Presenter
methods load rows, select customers, start saves, reset forms, and close
scopes. Commands are created there or by generated action companions consumed
there. The presenter should not build MigLayout panels. It should not subclass
\api{JPanel}. It may expose observable values, command sets, form models, and
row sources that the view binds.

The view-facing files should read like Swing. A panel creates components,
sets names or component keys, arranges them with MigLayout, installs generated
behavior plans, adapts commands to controls, and closes view-owned scopes. It
should not call repositories directly. It should not parse decimal text to
decide whether Save is enabled. It should not persist table state by itself
except through a controller whose store is supplied by the application. The
view's job is to present and bind.

Resource files should be reviewable by people who are not reading the Swing
code. Generated resource keys make this easier, but only if naming is
systematic. A resource map for the workspace can group command text, form
labels, table headers, tooltips, disabled reasons, validation messages, and
help hints. Tests can require the keys that must exist. Screenshot fixtures can
use the same resource map so the visible example shows the text that the
screen actually resolves.

Tests should follow the same ownership. A form-model test belongs near the
form model or generated source example. A presenter test belongs near the
presenter and uses fake services. A view-binding test runs on the EDT and
constructs the panel. A table-state test focuses on the descriptor, model,
table, and store. A screenshot test uses the screenshot harness and asserts
only enough structure to prove construction. One \api{CustomerWorkspaceTest}
file may contain several of these slices in a small example, but a large
application should let test files name the boundary they protect.

Screenshot generation should not be hidden in a manual recipe. A developer
should be able to find the Java source for a visual example, run the harness,
and understand which screen state is being captured. Captions and surrounding
text should mention the short example name, not a long filesystem path.

Avoid utility packages as dumping grounds. A helper that creates command
buttons from descriptors may belong in a Swing support package. A helper that
creates seeded customer rows may belong in test or example data. A helper that
loads resources may belong to application infrastructure. A file named
\api{SwingUtils} with every missing abstraction in it is a warning sign. The
same ownership rules apply to helpers: if the helper owns a durable fact,
name the fact; if it merely translates a boundary, keep it small.

The package shape should make generated code feel local rather than magical.
When \api{CustomerEdit} changes, the developer should know where to inspect
\api{CustomerEditFields}, \api{CustomerEditFormModel}, and
\api{CustomerEditBehaviorPlan}. When \api{CustomerRow} changes, the developer
should know where to inspect \api{CustomerRowColumns} and
\api{CustomerRowTableDescriptor}. When a presenter action changes, the
developer should know where to inspect the generated actions companion. The
build may place generated files under a generated-source directory, but the
conceptual ownership remains next to the annotated source.

\section{Review notes}

Review the screen by ownership. Which object owns row data? Which object owns
selection? Which object owns the form model? Which object owns commands? Which
object owns background tasks? Which object owns each binding and behavior?
Which object owns table state? Which object owns native dialog or window
disposal? If an answer is "the listener", look again. A listener should
translate an event; it should rarely own a durable fact.

This review habit is especially valuable when adding the eleventh feature. The
first ten features may each have seemed small enough for a local listener. The
eleventh exposes whether the screen has a structure. If adding validation
requires editing every text-field listener, the form model is missing. If
adding a menu item requires copying a toolbar listener, the command is
missing. If adding saved table state requires reverse-engineering column
indexes, the descriptor is missing. If closing the screen leaves a background
callback alive, the lifecycle owner is missing. The screen is maintainable
when these missing owners are easy to name and then easy to add.

Review the view by reading the layout constraints. Which region grows? Which
region has fixed width? Which row contains commands? Which area scrolls? Which
fields align by baseline? If the layout answer is hidden in nested panels and
magic preferred sizes, the screen will be hard to adapt. MigLayout makes the
answer compact, but the developer must still write the answer intentionally.

Review the command vocabulary. Are all user-visible operations represented by
commands or deliberate command-like models? Do keyboard shortcuts invoke the
same commands as buttons and menus? Do disabled reasons exist when users need
them? Are toggle states owned once? Does the command set expose the operations
tests and command palettes need? A screen with five ways to invoke Save should
have one Save command.

Review table identity. Are persisted column states keyed by stable ids? Are
view indexes converted at the boundary? Are renderers stateless? Is selection
preserved by domain identity rather than row number? Are generated or
handwritten descriptors the source of column metadata? The table is usually
the densest widget on a business screen, so table shortcuts taken early become
expensive later.

Review task boundaries. Which work can block? Where is cancellation recorded?
Where are stale results suppressed? Which values represent busy state,
progress, and failure? Which code returns to the EDT? Which lifecycle closes
the task handle? If a service callback directly mutates a component, the
screen has hidden threading policy.

Review screenshots last, but take them seriously. A screenshot can show
spacing, density, state presentation, tooltip text, busy overlays, and table
headers in a way prose cannot. It can also reveal that example references wrap,
captions look like body text, labels overflow, or a figure does not correspond
to the text. The screenshot is not decoration. It is a test artifact for the
book and a visual proof for the reader.

A final review pass should read the screen in the same order as a user. Start
at the top-left or first workflow entry point. Search or refresh. Select a row.
Move to detail. Edit a value. Observe validation. Save. Trigger help. Resize
the screen. Close it. At each step, ask which owner changes state and which
binding presents that state. This is more effective than reviewing files in
alphabetical order because it exposes missing links between otherwise correct
classes. A table descriptor can be perfect while selection still fails to load
detail. A form model can be correct while Ctrl+S bypasses active-editor
commit. A task service can be correct while the close path forgets to cancel
the current handle.

Then read the same screen in maintenance order. Rename a label. Add a column.
Add a command. Add a validation rule. Add a help topic. Add a screenshot state.
Change a service failure from modal dialog to inline banner. If each change
has an obvious owner, the architecture is healthy. If each change requires a
global search through listeners, the screen needs another named fact or
boundary. This maintenance-order review is particularly good for book
examples, because it asks whether the example teaches extensibility rather
than only initial construction.

Finally, check whether the screen can be explained in a few sentences without
lying. "Rows come from an observable list. The generated table descriptor
defines columns. Selection binding writes selected row state. A generated form
model edits the selected customer. Commands observe form, selection, and busy
state. Tasks deliver results to the EDT and stale results are suppressed.
Scopes close bindings and tasks. Screenshots render deterministic states."
If this explanation matches the code, the screen has earned its structure. If
one sentence requires an exception, go find the hidden owner before adding the
next feature.

The same explanation should survive visual polish. A designer may change
spacing, group fields differently, move commands into a toolbar, or choose a
more compact table header. Those changes should mostly touch the view and
resources. They should not require rewriting presenter workflow, table
identity, validation, or background-task policy. This is the practical test of
separation: visual improvement remains visual, workflow improvement remains
workflow, and stable identities keep the two from tearing each other apart.

It should also survive example growth. The compact book screen may later gain
filter chips, a column chooser, inline validation marks, a help pane, or a
second detail tab. Each addition should attach to an existing owner or justify
a new one. Filter chips attach to filter state and table descriptors. A column
chooser attaches to table state. Inline validation marks attach to problems.
A help pane attaches to help topics. A second detail tab attaches to a view
contract and form model boundary. Growth is healthy when new features make the
ownership map richer instead of bypassing it.

When that rule is followed, the example can stay small without being naive.
The book may show a compact workspace, while the reader can see where a larger
application would extend it. That is the right kind of simplification for a
technical example: fewer rows and services, not fewer owners.
The example remains honest because every omitted production feature has a
visible and reviewable place to attach when it becomes necessary.

\section{Failure patterns and recovery}

The first failure pattern is the listener-owned screen. It usually begins
innocently: a text field listener enables Save, a table selection listener
loads detail, a button listener starts Refresh, and a window listener cancels a
worker. Each listener is short. The screen works. Then validation must appear
in tooltips, command enablement, a problem summary, and tests. Selection must
survive refresh. Save must be available from a menu. The listener-owned screen
has no central facts to reuse. Recovery starts by naming facts, not by moving
listeners around. Introduce a form model, selected-row value, command, task
handle, and scope. Then let listeners translate Swing events into those facts.

The second failure pattern is string-owned identity. A table persists column
widths under visible headers. A test finds a button by English text. A help
router opens pages by tooltip text. A validation problem is known only by its
message. The screen looks clean until localization, copyediting, or
accessibility review begins. Recovery starts with stable keys. Add column
keys and table persistence ids. Add action keys and resource keys. Add problem
codes. Add help topics. Then migrate tests and persistence to those identities
before changing visible text.

The third failure pattern is table-as-application-state. A presenter reads
data from \api{JTable.getValueAt}, stores selection as view row number, and
updates cells by column indexes. This may be acceptable for a throwaway
prototype. It is fragile in a sorted, filtered, persistent, descriptor-backed
business screen. Recovery starts by defining a row type and row source. Then
install a table model from a descriptor, bind selection into presenter values,
and move command logic to row identity or selected row state. The table then
returns to its proper role: a presentation of row data.

The fourth failure pattern is command duplication. Save appears in a toolbar,
menu, keyboard shortcut, popup menu, and test helper. Each path calls a
slightly different method or checks a slightly different enabled condition.
Soon one path saves invalid data or ignores busy state. Recovery starts by
declaring one command identity and one command state. Existing buttons and
menu items can be adapted gradually. The shortcut can be rewired last. The
goal is not to use a command class for its own sake; the goal is that every
operation path observes the same operation model.

The fifth failure pattern is background work hidden in a component callback.
A button creates a worker. The worker captures the table. On success it
mutates rows directly. On failure it shows a dialog. On close it may or may
not be cancelled. This code is easy to write and hard to reason about.
Recovery starts by moving work into presenter/task state. The command starts
the task. The task reports through callbacks with an explicit EDT boundary.
Busy, progress, failure, and stale-result policy become observable facts. The
view binds to those facts. The worker no longer owns the screen.

The sixth failure pattern is overgeneration. A team sees generated form
companions and asks the processor to generate every panel, layout, service
call, permission check, renderer, and task. The result may remove some code,
but it also removes local judgment. Recovery starts by restoring the boundary.
Generate stable identity, descriptors, routine models, and routine bindings.
Keep layout, workflow, permissions, rendering, dynamic options, and failure
policy in handwritten code. The generated code should make the screen easier
to inspect, not harder to question.

The seventh failure pattern is screenshot drift. The application or its documentation has a
figure, but the figure is no longer generated from the example. The screen in
the figure shows different labels, spacing, or state from the current code. The
reader sees a polished image and learns an obsolete screen. Recovery starts by
making screenshot generation part of the example workflow. Use deterministic
data, fixed sizes, FlatLaf setup, and a generated figure path. Add at least a
smoke test that constructs the screenshot window. When the code changes, the
figure can be regenerated and reviewed instead of guessed.

The eighth failure pattern is cleanup by hope. A dialog closes, but document
listeners remain. A table-state controller still observes a table. A key
binding remains installed. A task callback holds the old view. The leak may
not appear in a short demo. It appears after repeated navigation or long
sessions. Recovery starts by listing installed relationships and assigning
them to scopes. Close scopes in the reverse direction of installation.
Generated behavior plans should participate in this ownership rather than
hide it.

The ninth failure pattern is premature performance folklore. A developer
avoids observable lists because "events are slow", replaces descriptors with
indexes because "objects are slow", or bypasses validation models because
"listeners are simpler". Sometimes performance work is needed. But in most
business screens, unclear ownership causes more slowness than named objects:
full table refreshes instead of precise events, repeated validation passes,
renderer state leaks, blocking EDT calls, and unnecessary repaint. Recovery
starts by measuring the actual slow path. Then optimize the owner that
controls it. Do not remove identity and lifecycle structure without evidence.

The tenth failure pattern is a screen that cannot be explained. Ask a
developer how Save becomes enabled, how a row selection loads detail, how a
stale refresh result is suppressed, how a field problem reaches the tooltip,
or how F1 chooses a topic. If the answer requires "let me search for the
listener", the screen may work but is not yet a teaching-quality screen.
Recovery is to write the explanation into the code through owners. This
chapter's architecture is not the only possible architecture, but any good
Swing workspace should be explainable in similarly concrete terms.

These patterns are not signs that Swing is obsolete. They are signs that the
screen has outgrown incidental code. Swing is quite capable of durable
business applications when meaning is not hidden in components. Corusco helps
by giving names and lifecycles to facts that otherwise become listener
fragments. The recovery technique is consistent: name the fact, choose the
owner, bind the presentation, test the boundary, and take a screenshot when
the result is visual.

\section{Exercises and checklist}

\begin{exercisebox}
\begin{enumerate}
\item Extend \bookAppExample{} with a generated-style component key for the
customer table and assert it in a test.
\item Replace the static table data with an \api{ObservableArrayList} and an
\api{ObservableTableModel}.
\item Add a selected-row value and bind table selection to it.
\item Add a Save command whose enabled state depends on dirty state,
validation, and busy state.
\item Add a simulated refresh task with stale-result suppression.
\item Add a generated or handwritten resource key for every visible label.
\item Add an F1 help topic for the name field and test the structured help
request.
\item Persist table column widths by stable ids, then change the visible
header text and verify that state restoration still works.
\item Close the view scope and assert that a later model change does not update
disposed components.
\item Generate a screenshot and review it with the local documentation output.
\end{enumerate}
\end{exercisebox}

\begin{center}
\textbf{Real-screen checklist.}
\end{center}

\begin{itemize}
\item Declare stable screen facts once.
\item Keep layout visible and reviewable.
\item Use MigLayout for ordinary production layout.
\item Use FlatLaf consistently for screenshots and examples.
\item Bind models to components through scopes.
\item Route buttons, menus, toolbars, and shortcuts through commands.
\item Convert table coordinates at explicit boundaries.
\item Persist table state by stable ids.
\item Model validation problems once and present them in several places.
\item Route background results through explicit EDT boundaries.
\item Suppress stale task results by generation or equivalent policy.
\item Close every binding, behavior, adapter, table controller, and task handle
with the lifecycle that installed it.
\item Test core, generation, Swing wiring, task behavior, and screenshot
construction separately.
\item Include screenshots for normal, busy, validation, table-state, or other
visual states that the chapter discusses.
\end{itemize}

A real screen is not made maintainable by one large abstraction. It is made
maintainable by many facts having the right owners. Records own committed
values. Form models own editing state. Descriptors own metadata. Commands own
operation state. Presenters own workflow. Views own components and layout.
Scopes own listener lifetimes. Tasks own background work. Resources own visible
text. Screenshots own visual evidence. Once those owners are visible, a Swing
workspace can grow without becoming a maze of incidental listeners.
